% !TEX root = ../linnik399.tex
\section{Assembly of the proof}\label{sec:join}

Each preceding argument holds once $q$ exceeds a threshold depending on fixed choices.
We first verify that these choices can be made in a noncircular order. A single
threshold then works for every leaf and every exterior case.

\subsection{The order of the choices}\label{subsec:order}
The constants of the argument are chosen in the following order. Each stage depends only on the
earlier ones, and there are finitely many choices at each stage.
\begin{enumerate}[(S1)]
  \item The fixed data: the exponent $L=3.99$ and the kernel of \cref{sec:detection}; the two far
    profiles of \cref{sec:far}; the test functions, anchors, sieve weights and grids of all near
    rows (\cref{sec:rows}); the case tree, its leaf programs and their certificates
    (\cref{sec:casetree,sec:lp}); the zero-location rows of \cref{sec:location}; the data of the
    exterior regimes (\cref{sec:exteriors}); $\eta=\eta'=10^{-6}$, $\lambda_{11}=0.05$ and $s_{\max}=3$; and
    $u_{\mathrm{exc}}=\min(1/\eta_{\mathrm{HB}}(\frac12),0.08)$, with $\eta_{\mathrm{HB}}(\delta)$ from \cref{inp:siegel}; and the envelope tolerance
$\eps_1=\eta/23$ of \cref{subsec:allowance}.
  \item The prime window: $C_P\ge\max(C_0(\eta H_0),3)$, with $C_0$ from \cref{inp:criterion}.
  \item The zero count $N=N(C_P)$ of \cref{lem:zerocount}; the mesh $h_{\mathcal A}$ and the finite set
    of anchors $\mathcal A$ of \cref{subsec:zerocount}; then the smoothing parameter $\theta$ of
    \cref{subsec:smoothing}, chosen so that \eqref{eq:thetachoice} holds for $\eps_1=\eta/23$; then the
    constant $c_{\mathrm{out}}$ of \cref{prop:firstoutside}.
  \item The count $K_0$ of \cref{subsec:representatives}, a bound, uniform in $q$, for the number
    of zeros of all nonprincipal $L$-functions modulo $q$ with $\lambda\le s_{\max}$ and $|\gamma|\le2$
    (\cref{inp:logfree}). Then the separation $M$, so large that for $|Y|\ge M$ we have
    $|\Re F(x+iY)|\le\eta''_{\min}/(2K_0+2)$ for every detector and $x\in[-s_{\max},0]$, where
    $\eta''_{\min}>0$ is the least of the numbers $\eta''=\eta'\min(1,\sqrt{I_BD_u},D_u)$ of
    \cref{cor:zeroform} over the finitely many rows, and
    $|\Re G(-\sigma+iY)|\le d_1$ for every Gram test and every $\sigma=s_1-\delta-\delta'$ that occurs
    (\cref{lem:decaytransform}). Then the buffer $C_B>\max(C_P+M,\,1.5)$ with
    $4c_{\mathrm{out}}/(H_0C_B^2)\le\eta$.
  \item The tolerances of the published results and of the lemmas proved here, each below the
    finitely many positive margins it must fit: the tolerances of \cref{inp:X31,inp:X32} for the
    finitely many test functions (those of \cref{sec:costs}, the functions $f^\theta_\sigma$ with tolerance
    $H_0\eps_1/8$, and the detectors of the near rows with tolerance $\eta''_{\min}/2$), the disc radius $\delta_0\le\frac13$ of \cref{subsec:zerocount},
    which is at most the radius that \cref{inp:X32} provides for each of these, and the tolerances
    of \eqref{eq:X34} and \cref{lem:explicitphi}, of
    \cref{inp:far} (with $\eps=\eta J^2$), of \cref{inp:burgess,inp:graham} and Mertens' estimate, of
    \cref{inp:criterion} (with $\eps=\eta H_0$), of the zero tables, and of the small-branch
    estimates of \cref{thm:small} (with $\mathcal E+\eps<0.03u_{\mathrm{exc}}$).
  \item Finally $q_0$: the maximum of the finitely many thresholds of the previous stages, together
    with those of \cref{inp:rectangle}, of the inclusions $R_P\subseteq R_B\subseteq$ the discs, and
    $\LL>4M(K_0+1)$ for \cref{lem:Tstar}.
\end{enumerate}
The certificates of stage (S1) depend only on fixed data, not on $C_P$, $M$, $C_B$ or $q$. The
count $K_0$ counts zeros up to height $2$, which exceeds $1+\delta_0$ for every admissible radius;
so $M$, which depends on $K_0$, can be fixed before the radius $\delta_0$, and there is no circular
dependence. The strip height $T^*=T^*(q)\in[\frac12,1]$ of \cref{lem:Tstar} is then determined by $q$.

Envelope errors that are summed over unboundedly many characters are weighted by
$e^{-A\lambda_\chi}$, and $\sum_\chi e^{-A\lambda_\chi}<19$ by \eqref{eq:Kfar}; errors in the near rows are
uniform per entry or are multiples of $(\sum_ja_j)^2$ (\cref{thm:graded}). Thus the error
allowances remain uniform as the number of characters grows with $q$.

\subsection{The regimes}
Let $q\ge q_0$. Exactly one of the following holds.
\begin{enumerate}[(R1)]
  \item No nonprincipal $L(s,\chi)$ has a zero in $R(l)$.
  \item $\lambda_1<0.1$.
  \item $0.1\le\lambda_1<1.5$.
  \item $\lambda_1\ge1.5$.
\end{enumerate}
(The first zero $\lambda_1$ is defined when (R1) fails; since $R(l)$ is compact there are finitely
many zeros in it.)

\subsection{Proof of \texorpdfstring{\cref{thm:main}}{the main theorem}}
Let $q\ge q_0$ and let $a$ be coprime to $q$.

In case (R1), $W=0$ and \cref{prop:detect} gives a prime $p\equiv a\pmod q$ with $q^A<p<q^L$.

In case (R2), if $\lambda_1<u_{\mathrm{exc}}$ then $P(a,q)\le q^{3.5}$ by \cref{prop:tiny}; if $u_{\mathrm{exc}}\le\lambda_1<0.1$ then
\cref{thm:small} gives a prime $p\equiv a\pmod q$ with $q^{3.6258}<p<q^{3.99}$.

In case (R3), \cref{thm:cover} gives a leaf of the case tree, or a subcase of a leaf, and a
feasible point of its program with $W\le\mathcal V(x)-2\eta<1-2\eta$. In case (R4), \cref{thm:large}
gives $W<1-2\eta$. In both cases \cref{prop:detect} gives a prime $p\equiv a\pmod q$ with
$q^A<p<q^L$.

Thus $P(a,q)<q^{3.99}$ for every $q\ge q_0$ and every $(a,q)=1$.
There are only finitely many moduli $2\le q<q_0$, and finitely many reduced residue
classes for each such modulus. Dirichlet's theorem supplies a least prime in each
class. Taking the maximum of $1$ and the finitely many ratios $P(a,q)/q^{3.99}$ gives
an absolute $C$ valid for all $q\ge2$. \qed

\begin{remark}[Effectivity]
Heath-Brown's theorem on Siegel zeros \cite[Corollary~1]{HeathBrown1990siegel}, used for
the closest exceptional zeros, is effective. The other source estimates used here are
also stated with effective constants. However, we have neither audited the effectivity
of every step of the assembly nor computed a value of $q_0$ or $C$.
\Cref{thm:main} asserts their existence and does not supply an explicit numerical bound
for either constant.
\end{remark}
